% !TEX program = pdflatex
\documentclass[11pt,a4paper]{article}
\usepackage[T1]{fontenc}
\usepackage[utf8]{inputenc}
\usepackage[brazil]{babel}
\usepackage{lmodern}
\usepackage{microtype}
\usepackage[margin=2.2cm]{geometry}
\usepackage{xcolor}
\usepackage{booktabs,longtable,tabularx,array}
\usepackage{enumitem}
\usepackage{titlesec}
\usepackage[hidelinks]{hyperref}

\definecolor{tocaInk}{HTML}{1E2127}
\definecolor{tocaHoney}{HTML}{C47E12}
\definecolor{tocaGrey}{HTML}{6B7280}
\titleformat{\section}{\large\bfseries\color{tocaInk}}{}{0pt}{}[{\color{tocaHoney}\titlerule[1pt]}]
\titlespacing*{\section}{0pt}{1.4em}{0.6em}
\setlist{itemsep=2pt,topsep=4pt}
\setlength{\parindent}{0pt}
\setlength{\parskip}{0.5em}

\newcommand{\disciplina}{{{TITLE}}}
\newcommand{\docente}{{{AUTHOR}}}

\begin{document}

{\color{tocaGrey}\small {{INSTITUTION}}}\par
{\LARGE\bfseries\color{tocaInk}\disciplina}\par
\vspace{0.3em}
{\color{tocaHoney}\rule{2.5cm}{3pt}}

\vspace{0.8em}
\begin{tabularx}{\linewidth}{@{}>{\bfseries}l X >{\bfseries}l X@{}}
  Docente & \docente & Semestre & {{YEAR}}.1 \\
  Carga horária & 60 h (4 créditos) & Horário & Terças e quintas, 10h--12h \\
  Contato & email@instituicao.br & Atendimento & Quartas, 14h--16h (sala 000) \\
  Pré-requisitos & Nenhum & Código & ECO0000 \\
\end{tabularx}

\section{Ementa}
Texto oficial da ementa: os grandes temas da disciplina em um parágrafo.

\section{Objetivos}
Ao final da disciplina, o estudante deverá ser capaz de:
\begin{enumerate}
  \item explicar ...;
  \item aplicar ...;
  \item avaliar criticamente ....
\end{enumerate}

\section{Conteúdo programático}
\begin{enumerate}[label=\textbf{Unidade \arabic*.},leftmargin=*,align=left]
  \item \textbf{Introdução.} Tópicos da primeira unidade.
  \item \textbf{Métodos.} Tópicos da segunda unidade.
  \item \textbf{Aplicações.} Tópicos da terceira unidade.
\end{enumerate}

\section{Cronograma}
\begin{longtable}{@{}c l p{0.5\linewidth} l@{}}
  \toprule
  Aula & Data & Tema & Leitura \\
  \midrule
  \endhead
  1 & 03/03 & Apresentação da disciplina & -- \\
  2 & 05/03 & Unidade 1: introdução & Cap. 1 \\
  3 & 10/03 & Unidade 1: continuação & Cap. 2 \\
  4 & 12/03 & Unidade 2: métodos & Cap. 3 \\
  5 & 17/03 & \textbf{Prova 1} & -- \\
  \bottomrule
\end{longtable}

\section{Avaliação}
\begin{tabular}{@{}lr@{}}
  Prova 1 & 30\% \\
  Prova 2 & 30\% \\
  Trabalho final & 30\% \\
  Participação & 10\% \\
  \midrule
  \textbf{Total} & \textbf{100\%} \\
\end{tabular}

Aprovação com média igual ou superior a 7,0 e frequência mínima de 75\%.

\section{Bibliografia}
\textbf{Básica}
\begin{itemize}
  \item SOBRENOME, Nome. \textit{Título do livro}. Cidade: Editora, ano.
\end{itemize}
\textbf{Complementar}
\begin{itemize}
  \item SOBRENOME, Nome. Título do artigo. \textit{Periódico}, v. 1, n. 1, p. 1--20, ano.
\end{itemize}

\end{document}
